\section{Three-dimensional exterior data and the mass inequality}\label{n3:intro:section}

\subsection{Initial data, normals, and enclosing area}

We work in three spatial dimensions, with zero cosmological constant and
\(G=c=1\). For a Riemannian metric \(g\) and a symmetric covariant
two-tensor \(K\), define the constraint densities by
\begin{equation}\label{n3:intro:constraints}
 16\pi\mu=R_g+(\tr_gK)^2-|K|_g^2,\qquad
 8\pi J=\Div_g\bigl(K-(\tr_gK)g\bigr).
\end{equation}
Here \(J\) is a covector field. The dominant energy condition is
\begin{equation}\label{n3:intro:dec}
                         \mu\ge |J|_g.
\end{equation}
All initial data considered here are smooth, including at a compact
boundary when one is present.

On an asymptotically flat end with Euclidean coordinates \(x\) and
\(r=|x|\), our basic assumptions are
\begin{equation}\label{n3:intro:decay}
 g_{ij}-\delta_{ij}=O_2(r^{-q}),\qquad
 K_{ij}=O_1(r^{-1-q}),\qquad q>\tfrac12.
\end{equation}
The notation \(O_j(r^{-a})\) includes the corresponding coordinate
derivative bounds through order \(j\). We assume that \(\mu\) and
\(|J|_g\) are integrable and that the following ADM limits exist and are
finite:
\begin{align}
 E&=\frac1{16\pi}\lim_{r\to\infty}
   \int_{S_r}(\partial_jg_{ij}-\partial_i g_{jj})n^i
                    \dd A_\delta,\label{n3:intro:energy}\\
 P_i&=\frac1{8\pi}\lim_{r\to\infty}
   \int_{S_r}\bigl(K_{ij}-(\tr_gK)g_{ij}\bigr)n^j
                    \dd A_\delta.\label{n3:intro:momentum}
\end{align}
The normal and area form in these definitions are Euclidean. Whenever
\(E>|P|_\delta\), write
\begin{equation}\label{n3:intro:rest-mass}
                         m=\sqrt{E^2-|P|_\delta^2}.
\end{equation}

For a two-sided surface \(\Sigma\) with specified unit normal \(\nu\),
we use
\[
 H_\Sigma=\Div_\Sigma\nu,\qquad
 \theta_+(\Sigma)=H_\Sigma+\tr_\Sigma K.
\]
Thus Euclidean spheres have positive mean curvature for the normal
toward infinity. A future marginally outer trapped surface, abbreviated
MOTS, satisfies \(\theta_+=0\). A weakly future outer trapped surface
satisfies \(\theta_+\le0\). Our spacetime sign convention is
\begin{equation}\label{n3:intro:k-sign}
 K(Y,Z)=\mathbf g(\mathbf\nabla_Y n,Z),
\end{equation}
where \(n\) is the future unit timelike normal. Equivalently, in Gaussian
normal coordinates the spatial metric has normal derivative \(2K\).

\begin{definition}[Exterior and enclosing area]\label{n3:intro:exterior-class}
An exterior is a connected orientable smooth three-manifold \(\Omega\)
with nonempty compact smooth boundary, complete as a metric space with
its boundary included, and with exactly one end, which is asymptotically flat.
An enclosing cut is a compact smooth embedded surface in
\(\overline\Omega\) separating the entire inner boundary from the
sufficiently distant part of that end. Equivalently, it bounds the side
containing the distant end; the other side is filled up to the inner
obstacle. A cut may have several components and may coincide with the
inner boundary. Its normal points toward the end. We put
\[
 a_g(\partial\Omega)=
  \inf\{|\Gamma|_g:\Gamma\text{ is an enclosing cut}\}.
\]
When using perimeter compactness, we take the closure of this class
among filled inner sets and include the area of any frontier coinciding
with the obstacle. Smooth outward approximation gives the same
infimum. Enclosing cuts and minimizing sets are taken with bounded
complementary pockets filled. Auxiliary perimeter arguments may use
competitors containing the inner obstacle before filling: filling a bounded
pocket deletes its frontier, cannot increase perimeter, and leaves the enclosing
infimum unchanged.
\end{definition}

The definition permits coincidence because first variations of the
minimum need not be realized by a cut lying strictly outside the
obstacle. It permits disconnected cuts because neither the minimizing
hull nor an intermediate trapped boundary need be connected. The
perimeter formulation and its compactness properties are established
in Section~\ref{n3:reduction:section}.

\subsection{The one-ended exterior inequality}

\begin{theorem}[Minimum-enclosing-area inequality]\label{n3:intro:exterior-inequality}
Let \((\Omega,g,K)\) be an exterior as in
Definition~\ref{n3:intro:exterior-class}. Assume
Equations~\eqref{n3:intro:constraints}--\eqref{n3:intro:decay}, integrability
of \(\mu\) and \(|J|_g\), and the finite ADM limits
\eqref{n3:intro:energy}--\eqref{n3:intro:momentum}. Orient \(\partial\Omega\)
by the normal into \(\Omega\), and suppose
\(\theta_+(\partial\Omega)\le0\). If \(E>|P|_\delta\), then
\begin{equation}\label{n3:intro:penrose}
          \sqrt{E^2-|P|_\delta^2}
               \ge \sqrt{\frac{a_g(\partial\Omega)}{16\pi}}.
\end{equation}
No extension of the data across \(\partial\Omega\) is required.
\end{theorem}
